\section{沿曲面切开：相空间是平坦联络}

\subsection{把时空切成两半}

量子化要从正则形式出发，而正则形式需要一个“时间”方向。三维流形 $M$ 没有天然的
时间，但我们可以沿一个闭曲面 $\Sigma$ 把它切开，写成 $\Sigma\times\RR$。
在 $\Sigma\times\RR$ 上，联络分解为
\[
A=A_0\,\dd t+A_i\,\dd x^i,
\]
其中 $A_0$ 是沿时间方向的联络分量（它是 $\mathfrak g$ 值的函数，不是一形式）。

把 (2.1) 代入并显式写出 $3$-形式的各项，注意 $\dd t\wedge\dd x^i\wedge\dd x^j$ 中
$i,j$ 只取两个空间方向，三次项只剩下 $A_0[A_i,A_j]$ 的组合，而 $A_0\wedge A_i\wedge A_j$
的符号使它与 $A_i\partial_tA_j$ 合并。取 $A_0=0$ 的规范，作用量变成
\begin{equation}
L=\frac{k}{8\pi}\int\dd t\int_\Sigma \varepsilon^{ij}\Tr\left(A_i\,\partial_tA_j\right)\dd^2x. \tag{3.1}
\end{equation}
这是一个对时间一阶的作用量：没有 $\dot A^2$ 项。它看起来像
$\int p\,\dot q$，其中 $p$ 就是 $A$ 的某个分量。因此相空间就是“所有 $A_i$”构成的空间，
并且它的辛形式可以立刻读出来。

\subsection{泊松括号}

从 (3.1) 读出正则对易关系。把 $\Sigma$ 上的两个分量记作 $A_1,A_2$（即 $A_i$ 的两维指标），
对 Lie 代数指标 $a,b$，有
\begin{equation}
\{A_i^a(x),A_j^b(y)\}=\frac{4\pi}{k}\,\varepsilon_{ij}\,\delta^{ab}\,\delta^2(x-y). \tag{3.2}
\end{equation}
注意括号里的 $4\pi/k$：耦合常数出现在辛形式的分母里。这正是为什么
“$k\to\infty$ 是弱耦合”在相空间里看得最清楚——$k$ 越大，$\hbar_{\text{eff}}=4\pi/k$ 越小。

\subsection{高斯定律与辛商}

$A_0$ 没有出现在 (3.1) 里，它是一个拉格朗日乘子。对 $A_0$ 变分给出约束：
\begin{equation}
\varepsilon^{ij}F_{ij}^a=0. \tag{3.3}
\end{equation}
在二维曲面上，曲率 $F_{ij}$ 只有一个独立分量（$F_{12}=-F_{21}$），所以 (3.3) 就是
\[
F=0,
\]
即联络必须平坦。这与第二讲从三维变分得到的结论一致：Chern--Simons 理论的经典解是平坦联络。

(3.3) 是“第一类约束”：它生成规范变换，因此不能简单地令其为零，而要按规范群取商。
用几何力学的语言，这是一个\textbf{辛商}（Marsden--Weinstein 约化）：先限制到约束面 $F=0$，
再按规范群的作用取商。两步合起来给出相空间
\begin{equation}
\cM=\left\{\text{$\Sigma$ 上的平坦联络}\right\}\big/\left\{\text{规范变换}\right\}
\cong\Hom(\pi_1\Sigma,G)/G. \tag{3.4}
\end{equation}
最后这个等式说明：平坦联络由它的和乐表示（$\pi_1\Sigma\to G$ 的同态）决定，只差一个共轭。

\subsection{相空间的维数}

$\Sigma$ 的亏格为 $g\ge2$ 时，$\pi_1\Sigma$ 有 $2g$ 个生成元，$\Hom(\pi_1\Sigma,G)$ 的维数是 $2gd$。
除以 $G$ 的共轭作用（$d$ 维），得到
\begin{equation}
\dim_\RR\cM=(2g-2)d,\qquad g\ge2. \tag{3.5}
\end{equation}
这个公式对 $g\ge2$ 成立。$g=1$ 时情形不同：一般位置的平坦联络有更大的稳定化子
（它的和乐落在某个极大环面里），所以不能简单地减去 $d$。对 $SU(2)$，环面上的模空间是
“枕头套” $T^2/\{\pm1\}$，实维数为 $2$。$g=0$ 时平坦联络只有平凡解，$\cM$ 是一个点。

(3.5) 是本书的第一个定量结论：相空间是有限维的，而且维数只依赖亏格与群，
不依赖任何连续参数。半经典地，量子 Hilbert 空间的维数应当约为
$\mathrm{Vol}(\cM)/(2\pi\hbar)^{n}$（$2n=\dim\cM$）；因为 $\cM$ 紧致、辛体积有限，这个数有限。
于是\textbf{Chern--Simons 理论在曲面上的 Hilbert 空间是有限维的}——这在三维量子场论里非常罕见，
也是整个故事能算得动的原因。

\subsection{复结构：从平坦联络到全纯丛}

要量子化，还需要给 $\cM$ 一个复结构。$\Sigma$ 上选一个复结构 $J$（等价地，一个亏格 $g$
的黎曼面结构）。Narasimhan--Seshadri 定理说：$\Sigma$ 上的平坦 $G$-联络与
“稳定的全纯 $G_\CC$-丛”之间存在自然的对应——一个平坦联络可以分解为
“全纯结构 + 与它相容的平坦幺正结构”。在这个对应下，$\cM$ 成为复流形 $\cM_J$，
而且它带有 Kähler 结构（对 $SU(N)$，$\cM_J$ 是 $c_1=0$ 的稳定秩 $N$ 全纯丛的模空间），
还是一个射影代数簇。

这个 Kähler 结构是下一讲量子化的入口：有了复结构，就能谈“全纯截面”，
而全纯截面正是我们要的量子态。

\subsection{辛商的具体形式：Atiyah--Bott}

约束面 $F=0$ 上的辛形式可以从 (3.2) 直接读出。设 $\delta A$ 是 $A$ 的一个变分，
对应切向量；辛形式是
\begin{equation}
\omega(\delta_1A,\delta_2A)=\frac{k}{4\pi}\int_\Sigma\Tr\left(\delta_1A\wedge\delta_2A\right). \tag{3.6}
\end{equation}
它显然是闭的（因为它是常系数二次型），并且在下述意义下非退化：若 $\delta_1A$ 与所有
满足 $F=0$ 的 $\delta_2A$ 的内积为零，则 $\delta_1A$ 必须是“纯规范”的。

Atiyah--Bott 的观察是：这个辛形式可以等价地写成模空间上的一个 Kähler 形式。
在 $\cM_J$ 上，切空间可以等同于 $H^1(\Sigma;\mathfrak g)$（一阶形变），
而 $\omega$ 就是 $H^1$ 上的 cup 积配对。于是 $\cM_J$ 的辛结构、复结构、
以及 (3.6) 的 Kähler 形式三者相容。这个等价性是第四讲能谈“Kähler 极化”的前提。

从 (3.6) 还可以看出 $k$ 的作用：$\omega$ 正比于 $k$。因此 $k$ 越大，辛体积越大，
半经典极限越好，量子 Hilbert 空间的维数越大——这与第四讲的 Weyl 律一致。

\subsection{小结}

把三维理论沿曲面切开之后，我们得到了一个具体得多的对象：有限维、紧致、Kähler 的模空间
$\cM_J$。三维量子场论的 Hilbert 空间，就是这个模空间上某个线丛的全纯截面空间。
剩下的问题是：哪个线丛？为什么是全纯截面？这是第四讲的主题。

\subsection*{练习}
\begin{enumerate}[leftmargin=*]
  \item 由 (3.1) 直接读出 (3.2) 的泊松括号（提示：把作用量写成 $\int p\dot q$ 的形式）。
  \item 计算 $g=0$ 与 $g=1$ 情形下 $\cM$ 的维数，并解释 $g=1$ 时公式 (3.5) 为什么不适用。
  \item 证明：$\cM$ 上由 (3.6) 定义的辛形式在下述意义下非退化——若它与所有闭形式的内积为零，则该形式是恰当的。
\end{enumerate}

\begin{quote}
\small
\textbf{本讲要点。} 正则形式给出辛形式与泊松括号 (3.2)；高斯定律 (3.3) 经辛商给出
平坦联络模空间 (3.4)，维数 $(2g-2)d$；选复结构后 $\cM_J$ 是 Kähler 模空间。
\end{quote}
